\documentclass[border=12pt]{standalone}
\usepackage[no-math]{fontspec}
\usepackage[brazil]{babel}
\usepackage{amsmath,amssymb}
\IfFileExists{C:/Windows/Fonts/Lato_Regular.ttf}{%
  \setmainfont{Lato}[
    Path = C:/Windows/Fonts/,
    UprightFont = Lato_Regular.ttf,
    BoldFont = Lato_Bold.ttf,
    ItalicFont = Lato_Italic.ttf,
    BoldItalicFont = Lato_BoldItalic.ttf
  ]
  \setsansfont{Lato}[
    Path = C:/Windows/Fonts/,
    UprightFont = Lato_Regular.ttf,
    BoldFont = Lato_Bold.ttf,
    ItalicFont = Lato_Italic.ttf,
    BoldItalicFont = Lato_BoldItalic.ttf
  ]
}{%
  \RequirePackage[default]{lato}
}
\renewcommand{\familydefault}{\sfdefault}

\usepackage{xcolor}
\usepackage{tikz}
\usepackage{pgfplots}
\pgfplotsset{compat=1.18}
\usetikzlibrary{arrows.meta,positioning,shapes.geometric,fit,backgrounds,calc,decorations.pathreplacing,patterns}

% Paleta institucional UEFS / Extensão Rural
\definecolor{erNavy}{HTML}{1F4268}
\definecolor{erNavyDark}{HTML}{152E4A}
\definecolor{erGold}{HTML}{B38A3D}
\definecolor{erGoldDark}{HTML}{8C6927}
\definecolor{erTeal}{HTML}{0E7C7B}
\definecolor{erForest}{HTML}{2D6A4F}
\definecolor{erAlert}{HTML}{9E2A2B}
\definecolor{erInk}{HTML}{26313D}
\definecolor{erMuted}{HTML}{5B6573}

\definecolor{bgNavy}{HTML}{EDF2F7}
\definecolor{bgGold}{HTML}{F8F3E9}
\definecolor{bgTeal}{HTML}{EBF6F5}
\definecolor{bgForest}{HTML}{EEF5F0}
\definecolor{bgAlert}{HTML}{FDF2F2}

\tikzset{
  >=Stealth,
  erBox/.style={
    draw=#1, fill=white, rounded corners=5pt, line width=1.15pt,
    inner sep=7pt, align=center, font=\footnotesize\sffamily, text=erInk
  },
  erTint/.style={
    draw=#1, fill=#1!8!white, rounded corners=5pt, line width=1.15pt,
    inner sep=7pt, align=center, font=\footnotesize\sffamily, text=erInk
  },
  erBadge/.style={
    fill=#1, text=white, rounded corners=9pt, inner xsep=8pt, inner ysep=3pt,
    font=\scriptsize\bfseries\sffamily
  },
  erTitle/.style={font=\normalsize\bfseries\sffamily, text=erNavy, align=center},
  erSub/.style={font=\footnotesize\sffamily, text=erMuted, align=center},
  erArrow/.style={->, >=Stealth, line width=1.35pt},
  erDArrow/.style={<->, >=Stealth, line width=1.35pt}
}

\hyphenpenalty=10000
\exhyphenpenalty=10000

\begin{document}
\begin{tikzpicture}[font=\sffamily]

  % Título principal
  \node[erTitle, font=\large\bfseries\sffamily] (title) at (0.65, 4.10)
    {Caixa de Ferramentas do Diagnóstico Participativo em Projetos de ATER};
  \node[erSub] (sub) at (0.65, 3.60)
    {Instrumentos metodológicos para leitura coletiva do território e planejamento rural};

  % Hub Esquerdo: Diagnóstico Participativo
  \node[
    draw=erNavyDark, fill=erNavy, rounded corners=8pt, line width=1.5pt,
    text=white, text width=3.35cm, minimum height=2.25cm, inner sep=10pt, align=center
  ] (hub) at (-6.05, 0.25) {%
    {\normalsize\bfseries Diagnóstico\\Participativo}\\[5pt]
    {\scriptsize Leitura coletiva\\da realidade rural}
  };

  % Grid 2x3 (Linha Superior: y = 1.85)
  \node[erTint=erTeal, text width=3.20cm, minimum height=1.70cm, inner sep=8pt] (g11) at (-1.35, 1.85) {%
    \textbf{\textcolor{erTeal}{DRP (Diagnóstico}}\\[1pt]
    \textbf{\textcolor{erTeal}{Rural Participativo)}}\\[4pt]
    \scriptsize\textcolor{erMuted}{Levantamento coletivo\\multitemático}
  };

  \node[erTint=erForest, text width=3.20cm, minimum height=1.70cm, inner sep=8pt] (g12) at (2.75, 1.85) {%
    \textbf{\textcolor{erForest}{Caminhada}}\\[1pt]
    \textbf{\textcolor{erForest}{Transversal}}\\[4pt]
    \scriptsize\textcolor{erMuted}{Observação direta das\\zonas agroecológicas}
  };

  \node[erTint=erNavy, text width=3.20cm, minimum height=1.70cm, inner sep=8pt] (g13) at (6.85, 1.85) {%
    \textbf{\textcolor{erNavy}{Mapa Falante da}}\\[1pt]
    \textbf{\textcolor{erNavy}{Propriedade}}\\[4pt]
    \scriptsize\textcolor{erMuted}{Uso do solo, água\\e infraestrutura local}
  };

  % Grid 2x3 (Linha Inferior: y = -1.35)
  \node[erTint=erGold, text width=3.20cm, minimum height=1.70cm, inner sep=8pt] (g21) at (-1.35, -1.35) {%
    \textbf{\textcolor{erGoldDark}{Linha do Tempo}}\\[1pt]
    \textbf{\textcolor{erGoldDark}{Histórica}}\\[4pt]
    \scriptsize\textcolor{erMuted}{Marcos produtivos e\\mudanças territoriais}
  };

  \node[erTint=erTeal, text width=3.20cm, minimum height=1.70cm, inner sep=8pt] (g22) at (2.75, -1.35) {%
    \textbf{\textcolor{erTeal}{Diagrama de Venn}}\\[1pt]
    \textbf{\textcolor{erTeal}{(Instituições)}}\\[4pt]
    \scriptsize\textcolor{erMuted}{Atores, parcerias e\\grau de influência}
  };

  \node[erTint=erNavy, text width=3.20cm, minimum height=1.70cm, inner sep=8pt] (g23) at (6.85, -1.35) {%
    \textbf{\textcolor{erNavy}{Matriz FOFA /}}\\[1pt]
    \textbf{\textcolor{erNavy}{SWOT}}\\[4pt]
    \scriptsize\textcolor{erMuted}{Forças, fraquezas,\\oportunidades e ameaças}
  };

  % Tronco central no corredor horizontal (y = 0.25) sem cruzamentos
  \draw[line width=1.45pt, color=erNavy] (hub.east) -- (6.85, 0.25);

  \draw[erArrow, color=erTeal]     (-1.35, 0.25) -- (g11.south);
  \draw[erArrow, color=erForest]   ( 2.75, 0.25) -- (g12.south);
  \draw[erArrow, color=erNavy]     ( 6.85, 0.25) -- (g13.south);

  \draw[erArrow, color=erGoldDark] (-1.35, 0.25) -- (g21.north);
  \draw[erArrow, color=erTeal]     ( 2.75, 0.25) -- (g22.north);
  \draw[erArrow, color=erNavy]     ( 6.85, 0.25) -- (g23.north);

\end{tikzpicture}
\end{document}
